We observe that the strong local performance of SOMs ultimately stems from the Newton step. Thus, to design accelerated methods with improved local performance, it is natural to seek second-order oracles that closely approximate the Newton step near the solution. One particularly promising candidate is the trust-region (TR) oracle~\cite{More1983,conn2000trust,yuan2015recent}. At a given point \( x \in \mathbb{R}^n \), the classic TR oracle solves the following subproblem:
\begin{equation}
    \label{eq.classic tr subproblem}
    \begin{aligned}
        \underset{d \in \mathbb{R}^n}{\min}~\tr(d) &:=     \nabla f(x)^T d + \tfrac{1}{2} d^T B(x) d \\
        \text{s.t.} \quad  \|d\|                          & \le \Delta,
    \end{aligned}
\end{equation}
where \( B(x) \) approximates the Hessian \( \nabla^2 f(x) \), and \( \Delta \) denotes the trust-region radius. For an overview of TR methods, we refer the readers to two excellent monographs \cite{yuan2000review,yuan2015recent}.

Two distinct strengths make TR oracle a compelling candidate for solving the fundamental question. On the theoretical side, \tr{} oracles can be solved in $O(\log\log(\epsilon^{-1}))$ arithmetic operations \cite{vavasisProvingPolynomialtimeSphereconstrained1990,yenewcomplexityresult1991}.  
Moreover, the ball constraint becomes (which can be detected by the multiplier being zero) near a non-degenerate minimizer, so the step reduces to the pure Newton direction, automatically revealing a region of quadratic convergence, which is an advantage absent from other second-order oracles. This unique feature offers a natural mechanism for local detection and exploitation. 
On the practical side, the TR oracle has been thoroughly analyzed \cite{jiangHolderianErrorBounds2022,rojasNewMatrixfreeAlgorithm2001}, so that highly competitive solvers have been developed \cite{rojasNewMatrixfreeAlgorithm2001, adachiSolvingTrustregionSubproblem2017, adachiEigenvaluebasedAlgorithmAnalysis2019} for it and its extensions.
To this end, TR forms the algorithmic core of general-purpose nonlinear programming solvers such as Knitro~\cite{byrd2006k}, COPT~\cite{geCardinalOptimizerCOPT2022}, and PDFO~\cite{ragonneau2024pdfo}. Similar success can be found in diverse applications ranging from adversarial training~\cite{yao2019trust} to reinforcement learning~\cite{schulman2015trust}. These advantages of \tr{} call for a systematic study of accelerating the TR method.

 Unfortunately, despite the strong empirical success of TR methods, the global complexity analyses remain incomplete. This gap has motivated several early efforts that focus on nonconvex problems~\cite{yeSecondOrderOptimization2005,curtis2017trust,curtis2021trust,zhangHomogeneousSecondorderDescent2025,hamad2022consistently,grapiglia2015convergence}. 
 % \cz{Here I added the 2015 MP paper, but bot by gratton. Check it.} 
While in convex optimization, other (non-accelerated) SOMs are known to achieve a convergence rate of \( O(\epsilon^{-1/2}) \) ~\cite{nesterov2006cubic,nesterov2008accelerating,mishchenko_regularized_2023}, the first TR method that matches this rate appeared only recently~\cite{jiang2024nonconvexityuniversaltrustregionmethod}. However, whether TR oracles can further benefit from acceleration, as achieved for other SOMs, remains an open question.  
% In this paper, we close this gap between TR methods and other accelerated SOMs. 

\paragraph{Challenges and our approach}
The proposal of using the TR oracle to address the global-local balance in accelerated SOMs introduces a central technical challenge: the Lagrange multiplier introduced by the TR constraint acts as a \textit{double-edged sword}.
\cz{I think it would be better to say the challenge of achieving local convergence for acc SOMs? because this paper is to do so}

\cz{This paragraph is a bit abstract.}
On one hand, as previously discussed, this multiplier captures valuable local geometric information and naturally vanishes near the optimal solution, allowing the TR oracle to closely recover pure Newton steps. To preserve this desirable local behavior, it is essential to exploit the multiplier's local feedback without imposing excessive intervention that could distort its geometric indication. In particular, maintaining the ability to approach Newton directions near the solution requires that the multiplier be allowed to adapt freely to the local landscape.

On the other hand, from the perspective of global convergence, it is often crucial to establish a stable relationship between the multiplier and the step size. This relationship underpins the ratio conditions commonly used in accelerated frameworks to control global progress. However, the multiplier, as a posterior quantity determined by the solution of the TR subproblem, tends to disrupt this ratio linkage, making it difficult to apply standard acceleration techniques and risking instability in the global complexity analysis.

This inherent tension reveals the core difficulty of our problem: achieving global acceleration requires regulating the multiplier’s influence to stabilize global progress, while preserving rapid local convergence demands minimal interference to retain its geometric sensitivity. Striking this balance calls for a careful refinement of the TR oracle to harmonize its global and local behaviors.

\cz{This part below seems to talk about what we do.}
To address this challenge, we build upon the recently proposed TR oracle of~\citet{jiang2024nonconvexityuniversaltrustregionmethod}, which stabilizes the relationship between the regularization and step size. However, directly applying this oracle is insufficient to resolve the tension between global stability and local precision, as the Lagrange multiplier can still destabilize the acceleration process if not carefully managed.

We introduce new strategies for adjusting the key parameters $\sigma$ (regularization strength) and $r$ (trust-region radius) in the TR oracle, which we call trust-region oracle for acceleration~(TROA):
\begin{subequations}\label{eq.newTR}
    \begin{align}
        \min_{d \in \mathbb{R}^n} \quad \trp(d) & := \nabla f(x)^T d + \tfrac{1}{2} d^T (\nabla^2 f(x) + \sigma I) d, \\\label{eq.sa.ballconstraint}
        \text{s.t.} \quad \|d\|                            & \le r.
    \end{align}
\end{subequations}
By carefully designing the adjustment of $\sigma$ and $r$, we allow the multiplier to retain its local geometric informativeness near optimality, which is crucial for recovering Newton-like behavior. At the same time, we ensure that the multiplier is globally stabilized through its interaction with the step size, which is essential for supporting acceleration.

In addition, we modify classical acceleration frameworks to accommodate the unique structure of the TR oracle. Specifically, we redesign the estimate sequence to handle multiplier-dependent behavior and ensure the validity of the acceleration dynamics under our modified oracle.

This combined strategy, which consists of adaptive parameter control and compatible acceleration design, enables us to balance the multiplier’s global and local influences and successfully overcome the core difficulty of applying acceleration to TR-type methods.

\paragraph{Contribution}
\cz{a bit too long and overlaps with the previous text.}
This work provides a constructive answer to the question posed earlier: we demonstrate, for the first time, that acceleration is achievable for second-order oracles while preserving their hallmark local efficiency. To this end, we extend the TR oracle proposed in~\citet{jiang2024nonconvexityuniversaltrustregionmethod} by incorporating greater flexibility in parameter selection, and develop two accelerated TR methods that reveal different aspects of the trade-off between global complexity and local convergence.

Our primary contribution is the \emph{Global-Local Accelerated Trust-Region Method (GL-ATR)} (Algorithm~\ref{alg.accelerated utr}), which directly targets the challenge of balancing global and local performance. GL-ATR is motivated by the observation that the Lagrange multiplier associated with the TR constraint encodes geometric proximity to optimality—its vanishing indicates that the algorithm is operating in a region where Newton-type behavior emerges. Leveraging this insight, we introduce a local detection mechanism (Algorithm~\ref{alg.local detection}) that dynamically activates acceleration when quadratic convergence regions are identified. This design yields an improved global worst-case oracle complexity of $\tilde{O}(\epsilon^{-1/3})$, while achieving a local convergence rate of $O(\log\log(1/\epsilon))$, consistent with quadratic convergence in high-accuracy regimes. To our knowledge, this is the first accelerated SOM that provably achieves both optimal global complexity and fast local convergence.

To further explore the global-local trade-off, we develop a second variant: the \emph{Near-Optimal Accelerated Trust-Region Method (NO-ATR)} (Algorithm~\ref{alg.ms accelerated utr}). NO-ATR is designed to push the limits of global performance, reaching the near-optimal global oracle complexity of $\tilde{O}(\epsilon^{-2/7})$ for convex problems. This is achieved by implementing an implicit search procedure tailored to the TR structure (Algorithm~\ref{alg.bisection ms}). However, our analysis and experiments reveal that this global optimality comes at the cost of local efficiency: unlike GL-ATR, NO-ATR does not retain superlinear convergence near the solution. This observation aligns with prior empirical findings~\cite{carmon2022optimal,huang2024inexact} and helps delineate the inherent trade-off between achieving global acceleration and preserving the fast local behavior of second-order methods.

% A comprehensive comparison of existing second-order accelerated methods and their complexity bounds is provided in Table~\ref{table.summary}. We remark that only our GL-ATR method exhibits strong local convergence performance among the methods in the table.

Our theoretical findings are corroborated by numerical experiments. On a global scale, both GL-ATR and NO-ATR perform comparably to the Accelerated Cubic Regularization method~\cite{nesterov2008accelerating} and the A-NPE method~\cite{monteiro2013accelerated}, while consistently outperforming the non-accelerated trust-region method~\cite{jiang2024nonconvexityuniversaltrustregionmethod} and the standard Cubic Newton method~\cite{nesterov2006cubic}. Locally, GL-ATR exhibits quadratic convergence, offering significant improvements over the local behavior of existing accelerated frameworks.

\cz{My version}
This work gives the first constructive proof that second-order oracles can be \emph{accelerated} without losing their Newton-type local speed. Building on the flexible trust-region (TR) oracle of \cite{jiang2024nonconvexityuniversaltrustregionmethod}, we design two algorithms that expose the global–local trade-off. The \emph{Global–Local Accelerated Trust-Region} method (GL-ATR) monitors the TR multiplier and triggers acceleration when it vanishes, achieving global oracle complexity $\tilde O(\epsilon^{-1/3})$ and local cost $O(\log\log\epsilon^{-1})$, consistent with quadratic convergence. The \emph{Near-Optimal ATR} variant (NO-ATR) embeds an implicit search tailored to the TR geometry, reaching the convex-case bound $\tilde O(\epsilon^{-2/7})$ at the price of forfeiting superlinear local speed. Numerical experiments confirm that both schemes match state-of-the-art global accelerators, while only GL-ATR retains quadratic convergence near the solution.

% \begin{table*}[t]
%     \caption{A Summary of the Mainstream Second-Order Oracles}
%     \label{table.summary}
%     \vskip 0.15in
%     \begin{center}
%         \begin{small}
%             \begin{sc}
%                 \begin{tabular}{lcccr}
%                     \toprule
%                     Oracle         & Nonconvex                 & Convex                                                  & Acceleration                     \\
%                     \midrule
%                     Cubic Newton      & $O(\epsilon^{-3/2})$~\cite{nesterov2006cubic} & $O(\epsilon^{-1/2})$\cite{nesterov2006cubic}                               & $O(\epsilon^{-1/3})$~\cite{nesterov2008accelerating}, $O(\epsilon^{-2/7})$~\cite{carmon2022optimal} \\
%                     GR Newton & $\tilde{O} (\epsilon^{-3/2})$~\cite{gratton2023yet}    & $O(\epsilon^{-1/2})$\cite{mishchenko_regularized_2023}                     & $\tilde O(\epsilon^{-1/3})$~\cite{doikov2020contracting}    \\
%                     Trust-region      & $O(\epsilon^{-3/2})^\dagger$                 & $O(\epsilon^{-1/2})$\cite{jiang2024nonconvexityuniversaltrustregionmethod} & $\tilde O(\epsilon^{-1/3})$, $\tilde O(\epsilon^{-2/7})$, this Paper                       \\
%                     \bottomrule
%                 \end{tabular}
%             \end{sc}
%         \end{small}
%     \end{center}
%     \vskip -0.1in
%     \begin{tablenotes}
%         \small
%         \item $\dagger$ denotes the works~\cite{curtis2017trust,curtis2021trust,luenberger_linear_2021,hamad2022consistently,hamad2024simplepracticaladaptivetrustregion,jiang2024nonconvexityuniversaltrustregionmethod}.
%     \end{tablenotes}
% \end{table*}

\subsection{Related works}
Beyond unconstrained minimization, a number of works~\cite{nesterov_lectures_2018, marques2022variants, dvurechensky2024improved, jiang2022generalized, lin2025perseus, huang2025approximation, ostroukhov2020tensor} have studied accelerated SOMs under stronger global assumptions. These include global strong convexity or similar uniform regularity conditions \cz{regularity conditions like global strong convexity}, under which they establish global linear, and local superlinear (or quadratic) convergence. \cz{what are \cite{marques2022variants, huang2025approximation} for? they are not linear? \cite{jiang2022generalized, lin2025perseus, ostroukhov2020tensor} are for min-max, which one is superlinear? the only superlinear result should be Dvurechensky's global superlinear results for self-concordant functions?} 

\cz{I think we should be more careful about the references we choose and present. Here is on my mind:
\begin{itemize}
\item talk about the \cite{nesterov_lectures_2018, marques2022variants, huang2025approximation}. A number of works have been proposed to study accelerated SOMs under stronger regularity conditions (e.g., global strongly convex functions), in which it is possible to achieve global linear...
More general settings have also been extensively discussed (min-max here)
\item In general, accelerated methods lose local superlinear convergence both theoretically and empirically. For example, \citet{jiang2024accelerated} proposed an accelerated quasi-Newton method that achieves the ..., while it underperforms the standard BFGS, which is known for superlinear local convergence.
\end{itemize}
}

\cz{somewhat duplicate text, we already mentioned this in the challenge and contribution}
In contrast, our work focuses on addressing the global-local trade-off in the general convex setting, without assuming global strong convexity. We only require the minimizer of the problem~\eqref{eq.main problem} to be non-degenerate, and our method is able to automatically detect when the iterates enter this local region without relying on restarts or prior knowledge of the local strong convexity parameters. To our knowledge, this is the first accelerated SOM that bridges the gap between global convergence and fast local rates under such minimal assumptions, making it applicable to a much broader class of convex problems.

\paragraph{Organization of the paper}
The remainder of the paper is organized as follows. Section~\ref{sec.priliminary} reviews the background on TR-type oracles and the use of estimate sequence techniques in second-order optimization. In Section~\ref{sec.first variant}, we present the GL-ATR, which achieves an oracle complexity of $\tilde O(\epsilon^{-1/3})$ and exhibits quadratic convergence in high-accuracy regimes. Section~\ref{sec.second variant} introduces the NO-ATR, which attains a near-optimal oracle complexity of $\tilde O(\epsilon^{-2/7})$. Section~\ref{sec.experiments} reports numerical results on both worst-case hard instances and practical tasks, confirming the theoretical advantages of our proposed methods.